\documentclass[11pt]{article}

\usepackage[a4paper,margin=19mm]{geometry}
\usepackage[T1]{fontenc}
\usepackage{lmodern}
\usepackage{microtype}
\usepackage{amsmath,amssymb,mathtools}
\usepackage{booktabs,array}
\usepackage{float}
\usepackage{siunitx}
\usepackage{xcolor}
\usepackage[hidelinks]{hyperref}

\definecolor{deepblue}{RGB}{0,79,128}
\definecolor{lightblue}{RGB}{238,246,251}
\definecolor{deeporange}{RGB}{205,92,20}
\definecolor{lightorange}{RGB}{254,244,235}

\sisetup{
  detect-all,
  retain-explicit-plus,
  table-number-alignment=center
}

\newcommand{\Ncal}{\mathcal{N}}
\newcommand{\proj}{\Pi_{[-1,1]}}

\title{One Projected-Gradient Descent Iteration\\
  for the Six-Bit, Two-Check Example}
\author{}
\date{}

\begin{document}
\maketitle
\vspace{-1.6em}

\begin{abstract}
This document lists every scalar used in one projected-gradient descent
iteration for the six-bit example in \texttt{p2}.  The calculation transforms
the six input log-likelihood ratios (LLRs) into an initial soft estimate,
evaluates the penalty function corresponding to the two parity-check
constraints, forms the complete gradient vector of this parity-only objective,
takes a gradient step, and projects the search point onto the box constraint.
\end{abstract}

\noindent
\fcolorbox{deeporange}{lightorange}{%
\begin{minipage}{0.945\textwidth}
\textbf{Scope.}  This exact projected-gradient teaching control, retained in
\texttt{p2}, uses the one-soft-estimate-per-coded-bit representation of the
project-defined CORAL receiver.  Its objective includes only the penalty terms
for the parity-check constraints: the part based on the received samples and
the penalty term for the binary constraint are omitted.  The step size
\(\eta=0.25\) is illustrative.  Production CORAL also
replaces the leave-one product on a violated check by its mapped value when
forming the search direction.  It is therefore not simply this calculation
with additional objective terms.
\end{minipage}}

\section{Graph, mapping, and initialization}

Bits are indexed from zero.  The parity-check matrix and graph neighbourhoods
are
\begin{equation}
H=
\begin{bmatrix}
0&1&1&0&0&1\\
1&0&0&1&1&0
\end{bmatrix},
\qquad
\Ncal(c_0)=\{1,2,5\},
\qquad
\Ncal(c_1)=\{0,3,4\}.
\label{eq:H}
\end{equation}
For a variable node, \(\Ncal(i)=\{c:i\in\Ncal(c)\}\) denotes its incident
checks.  The source convention maps bit one to \(+1\) and defines its supplied
input LLR by
\begin{equation}
x_i=2b_i-1,
\qquad
L_i=\log\frac{P(b_i=1\mid y_i)}{P(b_i=0\mid y_i)},
\qquad
x_i^{(0)}=\tanh\!\left(\frac{L_i}{2}\right).
\label{eq:init}
\end{equation}
Thus a positive input LLR gives a positive soft bit estimate.  The
hard-decision estimate at any iteration is
\begin{equation}
\widehat b_i^{(k)}=\mathbb{1}\!\left[x_i^{(k)}\geq0\right].
\label{eq:hard-decision}
\end{equation}
Because each
row of~\eqref{eq:H} contains three bits, a word that satisfies the corresponding
parity-check equation has bipolar product \(-1\).

\begin{table}[H]
\centering
\small
\caption{Observations, input LLRs, and initial soft bit estimates.}
\label{tab:init}
\begin{tabular}{
  c
  S[table-format=+1.2]
  S[table-format=+1.2]
  S[table-format=+1.3]
  S[table-format=+1.9]
  c
}
\toprule
{$i$} & {$y_i$} & {$L_i$} & {$L_i/2$} & {$x_i^{(0)}=\tanh(L_i/2)$}
& {$\widehat b_i^{(0)}$} \\
\midrule
0 & -0.24 & -0.38 & -0.190 & -0.187746206 & 0 \\
1 & +0.30 & +0.48 & +0.240 & +0.235495750 & 1 \\
2 & +0.09 & +0.14 & +0.070 & +0.069885890 & 1 \\
3 & +3.00 & +4.78 & +2.390 & +0.983347814 & 1 \\
4 & +1.79 & +2.85 & +1.425 & +0.890637366 & 1 \\
5 & +0.69 & +1.10 & +0.550 & +0.500520211 & 1 \\
\bottomrule
\end{tabular}
\end{table}

The \(y_i\) values identify the original observations.  This parity-only
calculation takes the supplied \(L_i\) values directly and does not derive them
from \(y_i\).

The displayed transmitted word is
\begin{equation}
\boldsymbol b_{\mathrm{tx}}=(0,1,0,1,1,1).
\label{eq:tx-word}
\end{equation}
The initial point and its hard-decision estimate are therefore
\begin{align}
\boldsymbol{x}^{(0)}={}&(-0.187746206,\ 0.235495750,\ 0.069885890,\nonumber\\
&\phantom{(}0.983347814,\ 0.890637366,\ 0.500520211),\\
\widehat{\boldsymbol b}^{(0)}={}&(0,1,1,1,1,1),
\qquad
H\widehat{\boldsymbol b}^{(0)}\bmod 2=(1,0).
\label{eq:initial-state}
\end{align}
Relative to \(\boldsymbol b_{\mathrm{tx}}\), only bit 2 is initially wrong, and
check \(c_0\) is violated.

\section{Objective and one-iteration rule}

For check \(c\), define
\begin{equation}
p_c(\boldsymbol{x})=\prod_{j\in\Ncal(c)}x_j,
\qquad
r_c(\boldsymbol{x})=1+p_c(\boldsymbol{x}),
\qquad
J_c(\boldsymbol{x})=r_c(\boldsymbol{x})^2.
\end{equation}
The plus sign is specific to the bit mapping and the degree-three checks in
this example.  For a degree-four check with the same mapping, a satisfying
word has product \(+1\), giving the corresponding squared residual
\((1-p_c)^2\).
The objective in this calculation is
\begin{equation}
J_p(\boldsymbol{x})=J_{c_0}(\boldsymbol{x})+J_{c_1}(\boldsymbol{x}).
\label{eq:objective}
\end{equation}
For \(i\in\Ncal(c)\), the partial derivative of \(J_c\) is
\begin{equation}
\frac{\partial J_c}{\partial x_i}
=2r_c(\boldsymbol{x})
 \prod_{j\in\Ncal(c)\setminus\{i\}}x_j.
\label{eq:check-derivative}
\end{equation}
The gradient of the objective function is obtained by adding the partial
derivatives from the checks that contain bit \(i\):
\begin{equation}
g_i^{(0)}=
\left.\frac{\partial J_p}{\partial x_i}\right|_{\boldsymbol{x}^{(0)}}
=\sum_{c\in\Ncal(i)}
 \left.\frac{\partial J_c}{\partial x_i}\right|_{\boldsymbol{x}^{(0)}}.
\label{eq:gradient}
\end{equation}
All six partial derivatives are evaluated at the same
\(\boldsymbol{x}^{(0)}\) before any variable is changed.  The Tanner graph
specifies which variables enter each product; these derivatives are not
directed variable-to-check or check-to-variable messages, and they are not
retained as state on directed edges.

One projected-gradient descent iteration consists of the gradient step and the
projection step
\begin{align}
z_i^{(0)}&=x_i^{(0)}-\eta g_i^{(0)},
&\eta&=0.25,
\label{eq:gradient-step}\\
x_i^{(1)}&=\proj\!\left(z_i^{(0)}\right)
=\min\!\left(1,\max(-1,z_i^{(0)})\right).
\label{eq:projection-step}
\end{align}
The Euclidean projection in~\eqref{eq:projection-step} enforces only the box
constraint \(\boldsymbol{x}\in[-1,1]^6\); it does not enforce parity.

\section{Objective at the initial point}

Table~\ref{tab:initial-checks} evaluates every quantity used by
\eqref{eq:check-derivative}.  The factor \(2r_c\) is included because it is
shared by the three partial derivatives associated with check \(c\).

\begin{table}[H]
\centering
\small
\caption{Initial check products, residuals, and penalty terms.}
\label{tab:initial-checks}
\begin{tabular}{c c
  S[table-format=+1.9]
  S[table-format=1.9]
  S[table-format=1.9]
  S[table-format=1.9]
  c}
\toprule
{check} & {bits} & {$p_c^{(0)}$} & {$r_c^{(0)}$} & {$2r_c^{(0)}$}
& {$J_c^{(0)}$} & {parity status} \\
\midrule
$c_0$ & $1,2,5$ & +0.008237477 & 1.008237477 & 2.016474953 & 1.016542809 & violated \\
$c_1$ & $0,3,4$ & -0.164429311 & 0.835570689 & 1.671141378 & 0.698178376 & satisfied \\
\bottomrule
\end{tabular}
\end{table}

Hence
\begin{equation}
J_p(\boldsymbol{x}^{(0)})
=1.016542809238+0.698178376170
\approx1.714721185407.
\label{eq:J0}
\end{equation}

\section{All partial derivatives at the initial point}

Define the product that excludes variable \(i\) from check \(c\) as
\begin{equation}
q_{c,i}^{(0)}=
\prod_{j\in\Ncal(c)\setminus\{i\}}x_j^{(0)}.
\end{equation}
Every variable belongs to one check in this example, so
\(g_i^{(0)}=\partial J_c/\partial x_i=2r_c^{(0)}q_{c,i}^{(0)}\).

\begin{table}[H]
\centering
\scriptsize
\caption{All products and partial derivatives used to form the initial gradient vector.}
\label{tab:gradient}
\setlength{\tabcolsep}{4.6pt}
\begin{tabular}{c c c
  S[table-format=+1.9]
  S[table-format=1.9]
  S[table-format=+1.9]
  S[table-format=+1.9]}
\toprule
{$i$} & {check} & {other bits} & {$q_{c,i}^{(0)}$} & {$2r_c^{(0)}$}
& {$g_i^{(0)}$} & {$-g_i^{(0)}$} \\
\midrule
0 & $c_1$ & $3,4$ & +0.875806306 & 1.671141378 & +1.463596157 & -1.463596157 \\
1 & $c_0$ & $2,5$ & +0.034979301 & 2.016474953 & +0.070534884 & -0.070534884 \\
2 & $c_0$ & $1,5$ & +0.117870382 & 2.016474953 & +0.237682674 & -0.237682674 \\
3 & $c_1$ & $0,4$ & -0.167213786 & 1.671141378 & -0.279437877 & +0.279437877 \\
4 & $c_1$ & $0,3$ & -0.184619821 & 1.671141378 & -0.308525822 & +0.308525822 \\
5 & $c_0$ & $1,2$ & +0.016457830 & 2.016474953 & +0.033186802 & -0.033186802 \\
\bottomrule
\end{tabular}
\end{table}

Thus
\begin{align}
\nabla J_p(\boldsymbol{x}^{(0)})={}&
(1.463596157,\ 0.070534884,\ 0.237682674,\nonumber\\
&\phantom{(}-0.279437877,\ -0.308525822,\ 0.033186802).
\label{eq:g0}
\end{align}

\section{Gradient step and projection step}

The change applied by the gradient step is
\(\Delta x_i^{(0)}=-\eta g_i^{(0)}\).  Table~\ref{tab:gradient-step}
lists every multiplication and addition in~\eqref{eq:gradient-step}.

\begin{table}[H]
\centering
\small
\caption{Gradient step with \(\eta=0.25\).}
\label{tab:gradient-step}
\setlength{\tabcolsep}{5.0pt}
\begin{tabular}{c
  S[table-format=+1.9]
  S[table-format=+1.9]
  S[table-format=+1.9]
  S[table-format=+1.9]}
\toprule
{$i$} & {$x_i^{(0)}$} & {$g_i^{(0)}$}
& {$\Delta x_i^{(0)}=-0.25g_i^{(0)}$} & {$z_i^{(0)}=x_i^{(0)}+\Delta x_i^{(0)}$} \\
\midrule
0 & -0.187746206 & +1.463596157 & -0.365899039 & -0.553645245 \\
1 & +0.235495750 & +0.070534884 & -0.017633721 & +0.217862029 \\
2 & +0.069885890 & +0.237682674 & -0.059420668 & +0.010465222 \\
3 & +0.983347814 & -0.279437877 & +0.069859469 & +1.053207283 \\
4 & +0.890637366 & -0.308525822 & +0.077131456 & +0.967768821 \\
5 & +0.500520211 & +0.033186802 & -0.008296701 & +0.492223511 \\
\bottomrule
\end{tabular}
\end{table}

Table~\ref{tab:projection} then applies the projection operator.  Only
coordinate 3 lies outside the box.

\begin{table}[H]
\centering
\small
\caption{Projection onto \([-1,1]^6\) and resulting hard-decision estimates.}
\label{tab:projection}
\begin{tabular}{c
  S[table-format=+1.9]
  S[table-format=+1.9]
  S[table-format=+1.9]
  c c}
\toprule
{$i$} & {$z_i^{(0)}$} & {$x_i^{(1)}$}
& {$x_i^{(1)}-z_i^{(0)}$} & {clipped?} & {$\widehat b_i^{(1)}$} \\
\midrule
0 & -0.553645245 & -0.553645245 & +0.000000000 & no  & 0 \\
1 & +0.217862029 & +0.217862029 & +0.000000000 & no  & 1 \\
2 & +0.010465222 & +0.010465222 & +0.000000000 & no  & 1 \\
3 & +1.053207283 & +1.000000000 & -0.053207283 & yes & 1 \\
4 & +0.967768821 & +0.967768821 & +0.000000000 & no  & 1 \\
5 & +0.492223511 & +0.492223511 & +0.000000000 & no  & 1 \\
\bottomrule
\end{tabular}
\end{table}

The point after one complete projected-gradient descent iteration is
\begin{align}
\boldsymbol{x}^{(1)}={}&(-0.553645245,\ 0.217862029,\ 0.010465222,\nonumber\\
&\phantom{(}1.000000000,\ 0.967768821,\ 0.492223511).
\label{eq:x1}
\end{align}

\section{State after one iteration}

Re-evaluating the two penalty terms at \(\boldsymbol{x}^{(1)}\) gives
Table~\ref{tab:post-checks}.

\begin{table}[H]
\centering
\small
\caption{Check values after the first projected-gradient descent iteration.}
\label{tab:post-checks}
\begin{tabular}{c c
  S[table-format=+1.9]
  S[table-format=1.9]
  S[table-format=1.9]
  S[table-format=1.9]
  c}
\toprule
{check} & {bits} & {$p_c^{(1)}$} & {$r_c^{(1)}$} & {$2r_c^{(1)}$}
& {$J_c^{(1)}$} & {parity status} \\
\midrule
$c_0$ & $1,2,5$ & +0.001122257 & 1.001122257 & 2.002244514 & 1.002245774 & violated \\
$c_1$ & $0,3,4$ & -0.535800606 & 0.464199394 & 0.928398787 & 0.215481077 & satisfied \\
\bottomrule
\end{tabular}
\end{table}

Therefore
\begin{align}
J_p(\boldsymbol{x}^{(1)})&=1.217726850659,\\
J_p(\boldsymbol{x}^{(0)})-J_p(\boldsymbol{x}^{(1)})
&=0.496994334749
\quad\text{(28.984\% relative to \(J_p(\boldsymbol{x}^{(0)})\)).}
\label{eq:cost-change}
\end{align}
The objective value decreases, but \(x_2^{(1)}=+0.010465222\) remains on the
positive side.  Consequently,
\begin{equation}
\widehat{\boldsymbol b}^{(1)}=(0,1,1,1,1,1),
\qquad
H\widehat{\boldsymbol b}^{(1)}\bmod2=(1,0).
\label{eq:final-hard}
\end{equation}
The hard-decision estimate and syndrome are unchanged: one iteration does not
yet correct bit 2.  In contrast, one sum-product message-passing iteration on
this same example changes bit 2 and recovers \(\boldsymbol b_{\mathrm{tx}}\).

Each variable node in this small graph has degree one, so the graph is a pair
of disconnected check stars.  It cannot demonstrate the aggregation of
several incident-check contributions into one gradient coordinate, and any
numerical similarity to sum-product does not establish equivalence.

\subsection*{Recomputed partial derivatives at \(\boldsymbol{x}^{(1)}\)}

These values describe the state available to the next iteration; they are not
used retroactively in the first update.

\begin{table}[H]
\centering
\scriptsize
\caption{Products and partial derivatives recomputed after the first iteration.}
\label{tab:post-gradient}
\setlength{\tabcolsep}{4.6pt}
\begin{tabular}{c c c
  S[table-format=+1.9]
  S[table-format=1.9]
  S[table-format=+1.9]
  S[table-format=+1.9]}
\toprule
{$i$} & {check} & {other bits} & {$q_{c,i}^{(1)}$} & {$2r_c^{(1)}$}
& {$g_i^{(1)}$} & {$-g_i^{(1)}$} \\
\midrule
0 & $c_1$ & $3,4$ & +0.967768821 & 0.928398787 & +0.898475400 & -0.898475400 \\
1 & $c_0$ & $2,5$ & +0.005151228 & 2.002244514 & +0.010314019 & -0.010314019 \\
2 & $c_0$ & $1,5$ & +0.107236813 & 2.002244514 & +0.214714320 & -0.214714320 \\
3 & $c_1$ & $0,4$ & -0.535800606 & 0.928398787 & -0.497436633 & +0.497436633 \\
4 & $c_1$ & $0,3$ & -0.553645245 & 0.928398787 & -0.514003574 & +0.514003574 \\
5 & $c_0$ & $1,2$ & +0.002279974 & 2.002244514 & +0.004565066 & -0.004565066 \\
\bottomrule
\end{tabular}
\end{table}

Thus the gradient ready for iteration two is
\begin{align}
\nabla J_p(\boldsymbol{x}^{(1)})={}&
(0.898475400,\ 0.010314019,\ 0.214714320,\nonumber\\
&\phantom{(}-0.497436633,\ -0.514003574,\ 0.004565066).
\end{align}

\paragraph{Numerical provenance.}
Within \texttt{JadeJcade}, the graph, observations, LLRs, and transmitted word
come from
\path{paper/p2/p2_assets/frames/os00_exact_example.tex}.  The objective and
illustrative step come from
\path{paper/p2/p2_assets/frames/os00b_projected_gradient.tex}; the retained
executable checker, \path{exact_six_bit_example.py}, is under
\path{paper/one_scalar_per_bit_slides/one_scalar_per_bit_slides_assets/examples/}.
Every displayed value was independently recomputed from the rounded input
LLRs.  A symmetric numerical derivative check with perturbation \(10^{-6}\)
agrees with
\(\nabla J_p(\boldsymbol{x}^{(0)})\) to within
\(1.92\times10^{-10}\) in every coordinate.

\end{document}
